% =================================================================
% Part XIV: Loom Final Architecture & Cumulative Findings
%   Frozen Behavioral Core + Sociable Shadow Layer
%   The empirical culmination of v8.3 → Loom v3.1
% =================================================================

\part{Loom Final Architecture and Cumulative Findings}
\label{part:loom_final}

\chapter{Loom v3.1 as Behavioral Core}
\label{ch:loom_v31_core}

\section{Position of Loom v3.1}

The empirical journey from v8.3 to Loom v3.1 culminated in a single
operational conclusion:

\begin{quote}
\emph{Loom v3.1 is the current Behavioral Core of the NRMO/Loom system.
It is not the maximum of any single world, but it is the most robust
across all worlds tested.}
\end{quote}

Loom v3.1 achieved \textbf{9/9 cells} in the Top~3 classification
across the canonical $3 \times 3$ benchmark (chaotic / drifting / noisy
$\times$ mild / moderate / severe). It is the first engine in the
NRMO research to do so without violating Sparse Activation
(\texttt{max\_active\_threads = 3}) and without weakening the v8.4.1
Safety Floor.

The architecture is summarised as:

\begin{equation}
\text{Loom v3.1}
  = \text{Safety Floor}_{\text{v8.4.1}}
  \;+\; \text{Specialist Threads}
  \;+\; \text{Sparse Activation}
  \;+\; \text{Oracle Gap Feedback}.
\end{equation}

Specialist threads are not external engines competing with Loom; they
are \emph{internalised} as named threads inside Loom v3.1. The
Behavioral Core decides, per step, which small number of threads to
fire.

\section{The Four Foundational Roles}

The Loom architecture stabilises around four irreducible roles:

\begin{description}
  \item[\textbf{v8.4.1 = Safety Floor.}]
    EmergencyResourceGuard, ActionIntensityThrottle, Revalidation,
    and CumulativeRiskTracker. \emph{Immutable.} Removing or relaxing
    any of these reproduces the v8.3 aggressive-overdrive failure
    mode. Safety Floor is never bypassed except inside the strictly
    bounded ContextualMerger-bypass under Hard Drift mode, where it
    is still applied as a final validation step.
  \item[\textbf{v9\_minimal = Drift Specialist.}]
    A minimal-intervention engine that dominates DriftingWorld by
    \emph{not interfering with regime shift}. Inside Loom, this
    behaviour is absorbed as the \texttt{DriftThread}, fired
    sparingly when drift is detected.
  \item[\textbf{v8.5.1 = Stabilization Specialist.}]
    Strong in chaotic/noisy mild--moderate cells. Inside Loom, this
    behaviour is absorbed as the \texttt{StabilizationThread}, using
    Recovery + Defensive + Synthesis under ContextualMerger.
  \item[\textbf{ActiveCycle = Severe/Volatile Specialist.}]
    Provides cycle-aware adaptive behaviour for the high-volatility
    regime. Inside Loom, this is the \texttt{SevereCycleThread}.
\end{description}

\section{Why Loom v3.1 is the Correct Frozen Baseline}

A natural objection is: \emph{Why not promote v9\_minimal as the
canonical engine, since it dominates DriftingWorld?}

The answer is that v9\_minimal is \emph{not} globally optimal. Its
strengths and weaknesses are:

\begin{itemize}
  \item \textbf{Strength.} Minimal correction allows it to follow
    drift trajectories cleanly. It does not over-defend, does not
    over-recover, and does not over-merge.
  \item \textbf{Weakness.} Its Safety Floor is thin. In chaotic/noisy
    regimes the Stabilization-style behaviour outperforms it. Its
    lower-tail and ruin control require external validation.
\end{itemize}

NRMO's purpose is not to maximise score in a single world. NRMO's
purpose is to maintain reachability of admissible futures while
respecting the ruin boundary across all worlds. Loom v3.1 satisfies
this constraint; v9\_minimal does not.

Therefore:

\begin{quote}
\emph{v9\_minimal is the Drift Specialist to be absorbed.\\
Loom v3.1 is the Behavioral Core.}
\end{quote}

\chapter{Sociable Shadow Layer}
\label{ch:sociable_shadow}

\section{The Theory Update: Detection $\neq$ Intervention}

The single most important theoretical contribution of this
research cycle is the following principle, validated empirically
against the v3.2 / v3.2.1 negative results:

\begin{theorem}[Detection--Intervention Separation]
\label{thm:detection_neq_intervention}
In a Loom-style control architecture with a Safety Floor and Sparse
Activation, detection signals (cycle/orbit/failure-face/canonical
state) may be acquired continuously without harm, but routing those
signals into immediate per-step action interventions destabilises
the controller and degrades score.
\end{theorem}

This is not a defect of the underlying sociable-numbers theory. The
detection mechanisms themselves --- orbit canonicalisation, cycle
recurrence scoring, drift-escape estimation, failure-face
classification --- behave exactly as designed. The error in v3.2 and
v3.2.1 was \emph{control-theoretic}, not detection-theoretic: a
high-frequency detector was wired directly into a mode switch
without hysteresis, dwell time, or confidence gating, producing
over-control, mode hunting, and loss of Drift-mode purity.

The corrected design assigns the sociable-numbers machinery to a
\emph{shadow} layer.

\section{Architecture of the Sociable Shadow Layer}

The Sociable Shadow consists of four observational components,
all default ON:

\begin{description}
  \item[\textbf{SociableObservationLayer.}]
    Records every step as $(state, action, module, context, world,
    reward, \Delta, \mathrm{guard}, \mathrm{throttle}, \mathrm{reval})$
    with bucketed signatures.
  \item[\textbf{SociableCanonicalizationLayer.}]
    Assigns canonical IDs to state, failure, context, and orbit.
    Same-orbit, same-failure, and same-context structures are
    collapsed into canonical equivalence classes.
  \item[\textbf{SociableCycleOrbitDetector.}]
    Computes orbit-recurrence, cycle-closure, drift-escape,
    directional-persistence, reversal-rate, autocorrelation, and
    no-improvement-cycle scores. From these it derives
    likelihood scores for drifting, chaotic, and noisy world types.
  \item[\textbf{SociableFailureFaceDetector.}]
    Classifies failures into a small enumerated set of faces
    (\texttt{R\_CRITICAL}, \texttt{X\_HIGH}, \texttt{GUARD\_FORCED},
    \texttt{THROTTLE\_FORCED}, \texttt{STABILIZATION\_OVERUSE},
    \texttt{DRIFT\_MISS}, \texttt{RECOVERY\_DOMINANCE},
    \texttt{AGGRESSIVE\_DRAWDOWN}, \texttt{OVER\_DEFENSE},
    \texttt{NO\_IMPROVEMENT\_CYCLE}, \dots) and tracks their
    distribution and entropy.
\end{description}

All four are written into the decision trace. None of them write back
into Loom v3.1's action selection. Empirically, this gives a
\emph{paired diff of exactly zero} between Loom v3.1 and Loom v3.1 +
Shadow across all nine cells (Wilcoxon test: all-zero differences,
$p$ undefined / trivial).

\section{Default-ON vs Default-OFF Components}

\begin{center}
\begin{tabular}{lll}
\toprule
\textbf{Component} & \textbf{Default} & \textbf{Reason} \\
\midrule
SociableObservationLayer & ON & observation only \\
SociableCanonicalizationLayer & ON & passive labelling \\
SociableCycleOrbitDetector & ON & metric computation only \\
SociableFailureFaceDetector & ON & classification only \\
EffectiveRangeGuard & ON & coverage bookkeeping \\
\midrule
FF-driven thread suppression & OFF & changes action \\
Hard/Soft Drift mode override & OFF & changes mode \\
Cycle-driven action suppression & OFF & changes action \\
FF-driven Merger bypass & OFF & changes selection \\
\bottomrule
\end{tabular}
\end{center}

\noindent The OFF components remain available for future use, but
require additional control structure (hysteresis, dwell time,
confidence threshold, oracle-gap validation) before they can be
re-enabled without reproducing the v3.2 regression.

\chapter{The v3.2 / v3.2.1 Negative Results}
\label{ch:v32_negative}

\section{What Was Attempted}

Loom v3.2 attempted to wire the Sociable Detection layers directly
into Sparse Activation. Loom v3.2.1 reduced the Hard/Soft Drift
thresholds and tightened SevereCycle conditions. Both versions
verified that:

\begin{itemize}
  \item Drift mode activation rose dramatically (HardDrift: $1\,320$,
    SoftDrift: $1\,483$ in drifting/mild per 100 runs).
  \item Failure faces were classified correctly
    (\texttt{r\_critical}, \texttt{guard\_forced}, \texttt{x\_high},
    \texttt{drift\_miss}, \dots).
  \item ContextualMerger bypass was triggered as designed.
  \item All four detection layers operated to specification.
\end{itemize}

\section{What Went Wrong}

Despite correct detection, scores degraded:

\begin{center}
\begin{tabular}{lrrr}
\toprule
\textbf{Cell} & \textbf{v3.1} & \textbf{v3.2} & \textbf{v3.2.1} \\
\midrule
drifting/mild     & 19.34 & 11.64 & 10.75 \\
drifting/moderate &  9.93 &  8.91 &  7.70 \\
drifting/severe   &  6.40 &  5.96 &  6.09 \\
\midrule
Top-3 count       & 8--9/9 & 4/9   & 2/9 \\
\bottomrule
\end{tabular}
\end{center}

The mechanism of degradation was:

\begin{enumerate}
  \item Detection signals fluctuated step by step.
  \item Each fluctuation caused a mode switch (Drift $\leftrightarrow$
    Stabilization $\leftrightarrow$ FF-intervention).
  \item Frequent switching destroyed the purity of Drift mode, which
    requires \emph{not interfering} with trend.
  \item ContextualMerger bypass became too frequent, replacing
    well-merged candidates with raw V71 base proposals that were no
    longer the best choice in the current context.
  \item The Behavioral Core lost coherence and underperformed even
    its own pre-detection baseline.
\end{enumerate}

\section{The Lesson}

\begin{quote}
\emph{Increasing detection precision does not increase control
performance. Feeding detection signals directly into control without
hysteresis, dwell time, or confidence gating reproduces a classical
hunting / over-control instability.}
\end{quote}

This is the principle stated formally in
Theorem~\ref{thm:detection_neq_intervention}. v3.2 and v3.2.1 are
preserved as the negative-result evidence supporting it.

\chapter{Specialist Integration: Final Map}
\label{ch:specialist_integration_final}

\section{The Four-Role Diagram}

\begin{center}
\fbox{\parbox{0.85\textwidth}{
\centering
\textsc{Loom Final Architecture (current frozen state)} \\[1ex]
\begin{tabular}{l@{\hspace{1.5em}}l}
v8.4.1            & Safety Floor (immutable)                       \\
v9\_minimal       & Drift Specialist (absorbed as DriftThread)     \\
v8.5.1            & Stabilization Specialist (absorbed)            \\
ActiveCycle       & Severe / Volatile Specialist (absorbed)        \\
\midrule
Loom v3.1         & \textbf{Behavioral Core (frozen)}              \\
Sociable Shadow   & Detection / Trace Layer (observation only)     \\
\end{tabular}
}}
\end{center}

\section{Why This Map Is Stable}

Each role is necessary, each role is non-overlapping, and each role
has been individually validated by ablation:

\begin{itemize}
  \item Removing the Safety Floor reproduces v8.3.
  \item Replacing the Behavioral Core with any single Specialist
    forfeits one or more worlds.
  \item Re-enabling Sociable Intervention reproduces v3.2 / v3.2.1.
  \item Removing the Sociable Shadow removes the trace layer that
    enables offline analysis, but does not affect runtime behaviour.
\end{itemize}

Therefore the configuration

\begin{equation}
\text{NRMO}_{\text{current}}
  = \underbrace{\text{Loom v3.1}}_{\text{behaviour}}
  \;\;+\;\;
  \underbrace{\text{Sociable Shadow}}_{\text{observation}}
\end{equation}

is the empirically minimal stable configuration that simultaneously
preserves the Safety Floor, the Specialist behaviours, Sparse
Activation, and offline interpretability.

\chapter{Cumulative Findings}
\label{ch:cumulative_findings}

The following ten findings are the durable empirical results of the
v8.3 $\rightarrow$ Loom v3.1 + Shadow research cycle. They are
intended to outlive any individual implementation.

\begin{enumerate}
  \item \textbf{No universal best engine exists.} No single engine
    won every cell of the $3 \times 3$ benchmark. Empirical fact.
  \item \textbf{Specialism dominates Generalism per world.} Each
    world had a specialist that outperformed every generalist
    in that world.
  \item \textbf{v8.4.1 Safety Floor is immutable.} Every relaxation
    attempt re-introduced v8.3-class failure modes.
  \item \textbf{Synthesis Pathway is universally positive.} The
    only sub-module that contributed non-negatively in every cell.
  \item \textbf{The Hard Guard is EmergencyResourceGuard.} The other
    components of v8.4.1's safety apparatus are supportive; the
    Emergency Guard alone carries the protective load.
  \item \textbf{Sparse Activation works.} Limiting
    \texttt{max\_active\_threads} to 3 improves stability without
    sacrificing capability.
  \item \textbf{Drift Override + Enhanced Drift Thread is the
    correct drifting strategy.} Detection alone is not enough;
    minimal-intervention execution is required.
  \item \textbf{Sociable Detection belongs in trace, not control.}
    Verified by the v3.2 / v3.2.1 negative results.
  \item \textbf{Excessive mode switching degrades score.} A control
    architecture must have dwell time, not just precision.
  \item \textbf{Detection $\neq$ Intervention.} The Loom theory
    update. Observation may be continuous; intervention must be
    gated.
\end{enumerate}

\chapter{One-Line Summary}
\label{ch:one_line_summary}

\begin{quote}\large
\emph{Loom v3.1 is the current behavioral core.
v9\_minimal is the Drift Specialist to be absorbed, not the final
engine. v8.4.1 remains the immutable Safety Floor. Sociable Essence
remains always-on as Shadow Detection, never as direct
intervention.}
\end{quote}

\bigskip
\noindent
\textsc{Final architectural identity:}
\begin{equation*}
\boxed{\;
  \text{NRMO}_{\text{current}}
  \;=\;
  \text{Loom v3.1 (frozen behavioral core)}
  \;+\;
  \text{Sociable Shadow Layer (observation only)}
\;}
\end{equation*}
